\documentclass[10pt,a4paper]{amsart}
\usepackage[margin=19mm]{geometry}
\usepackage{amsmath,amssymb,mathtools}
\usepackage[scr=rsfs,cal=euler]{mathalfa}
\usepackage{xfrac}
\usepackage[unicode,hidelinks,bookmarks=false]{hyperref}
\newcommand{\ie}{i.e.\ }
\newcommand{\cf}{cf.\ }
\newcommand{\resp}{respectively\ }
\newcommand{\Subsection}{\subsection}
\providecommand{\nobreakdash}{\nobreak\hbox{-}}
\setlength{\emergencystretch}{2em}
\multlinegap0pt
\usepackage{bm}
\usepackage{bbm}
\usepackage{dsfont}
\usepackage{xspace}
\usepackage{cancel}
\usepackage{mathtools}
\usepackage{amsthm}
\makeatletter
\@ifundefined{theorem}{\newtheorem{theorem}{Theorem}}{}
\@ifundefined{lemma}{\newtheorem{lemma}[theorem]{Lemma}}{}
\@ifundefined{proposition}{\newtheorem{proposition}[theorem]{Proposition}}{}
\@ifundefined{remark}{\newtheorem{remark}[theorem]{Remark}}{}
\makeatother
\newcommand{\C}{\mathbf{C}}
\title[Normal forms and geometric structures on Hopf manifolds]{Normal forms and geometric structures on Hopf manifolds}
\author{Paul Boureau}
\date{Source: arXiv:2501.10346v2 (CC BY 4.0)}
\begin{document}
\maketitle

\noindent\emph{Source and licence.} Normal forms and geometric structures on Hopf manifolds. Version arXiv:2501.10346v2 (CC BY 4.0), distributed under the Creative Commons Attribution 4.0 International licence, as recorded in the arXiv licence metadata for this exact version and shown on the abstract page https://arxiv.org/abs/2501.10346v2. Full rights notice: licenses/arxiv-2501.10346-CC-BY-4.0.txt.\par\smallskip

\noindent\textit{Editorial scope.} Counted unit: the Proposition whose source label is \texttt{p:Ggroup} in the article (source0.tex lines 514--516, source label \texttt{p:Ggroup}), delivered together with the article's own complete original proof of it (source0.tex lines 518--565, one \texttt{proof} environment of 48 lines). No mathematics is written, repaired or re-derived: statement and proof are the article's own text line for line. Required context retained uncounted: r:monomial (lines 251--256). Every definition, display and label of those lines is retained unchanged. Omitted from this package and not redistributed: 1--250, 257--513, 517--517, 566--1028. Nothing inside a retained excerpt is omitted or altered: every retained body is delivered exactly as the article's lines are written, so no omission receipt of any kind accompanies this package. Citations: the retained excerpts carry no citation command at all, so no bibliography entry of the article is transcribed and no reference is redistributed. Reference adaptation: none was needed and none was made; every reference inside the delivered unit resolves inside it, so no unresolved reference or citation is produced. Printed slips kept literal: no printed word, symbol, hypothesis or number of the counted statement or of its proof was altered. Layout and class: the article's own class is \texttt{amsart} with packages including amssymb, amsmath, amsthm, srcltx, babel, graphicx, color; the wrapper is the lane's pinned \texttt{amsart} editorial wrapper at 10pt on A4, which restates the theorem-like environments the retained text needs and adds the article's own macro definitions that the retained text needs (\texttt{\textbackslash C}), with extra wrapper packages (bm, bbm, dsfont, xspace, cancel, mathtools). The delivered numbering is the unit's own. Assets: no retained span contains a figure environment, a graphics-inclusion command, a PDF-page inclusion, a table or a TikZ picture; no image file is copied into this package, the package declares no asset dependency and no image file is redistributed with it. Licence: version arXiv:2501.10346v2 of this article is distributed under the Creative Commons Attribution 4.0 International licence, as recorded in the arXiv licence metadata for this exact version and shown on the abstract page https://arxiv.org/abs/2501.10346v2. A full rights notice is delivered at licenses/arxiv-2501.10346-CC-BY-4.0.txt. Characters: every retained excerpt is pure ASCII, so the delivered engine, XeLaTeX with its default Latin Modern text font, reports no missing character.
\begin{remark}[monomial criterion]\label{r:monomial}
In the adapted coordinates, write $z^I:=z_1^{i_1}\cdots z_n^{i_n}$ for a multi-index $I=(i_1,\dots,i_n)$, and $H_{I,j}:=z^Ie_j$. For $1\le k\le n$ let $b(k)\in\{1,\dots,l\}$ be the block of the coordinate $k$, so that $|\lambda_k|=\mu_{b(k)}$. The monomial map $H_{I,j}$ is homogeneous of type $s(I)$, where $s_b(I):=\sum_{k:\,b(k)=b} i_k$, and it is sub-resonant if and only if
\[
\ln\mu_{b(j)}\ \le\ \sum_{b}s_b(I)\ln\mu_b\ =\ \sum_{k} i_k\ln|\lambda_k| .
\]
\end{remark}

\begin{proposition}\label{p:Ggroup}
$G := \langle SR^*(L), \C^n \rangle$, where $\C^n$ acts by translation, is an algebraic group of finite dimension.
\end{proposition}

\begin{proof}
For $\tau \in \C^n$, let $t_\tau$ denote the map $z \mapsto z + \tau$. We begin by proving
\begin{lemma}\label{l:Gform}
\[
G = \{t_\tau \circ h \mid \tau \in \C^n, h \in SR^*(L)\}.
\]
\end{lemma}
\begin{proof}
We need to verify that for every $(\tau, h) \in \C^n \times SR^*(L)$,
\[
t_{-h(\tau)} \circ h \circ t_\tau \in SR^*(L).
\]

The group structure of $SR^*(L)$ ensures the invertibility of the derivative at $0$. Let us demonstrate the sub-resonant nature of the result. To this end, note that for a fixed $\tau$, the map $h\mapsto t_{-h(\tau)}\circ h\circ t_\tau$, that is $h\mapsto h(\cdot+\tau)-h(\tau)$, is linear in $h$. Since in the adapted coordinates every element of $SR(L)$ is a linear combination of sub-resonant monomials $H_{I,j}$ (Remark \ref{r:monomial}), it suffices to verify the condition for such monomials.

We have:
\[
t_{-H_{I,j}(\tau)} \circ H_{I,j} \circ t_\tau (z)
= \big((z_1 + \tau_1)^{i_1} \cdots (z_n + \tau_n)^{i_n} - \tau_1^{i_1} \cdots \tau_n^{i_n}\big)\, e_j .
\]
Expanding, the homogeneous terms of $t_{-H_{I,j}(\tau)} \circ H_{I,j} \circ t_\tau$ are scalar multiples of the monomials $z^J e_j$ with $J\le I$ componentwise and $J\neq0$. Each of them is sub-resonant: by Remark \ref{r:monomial} and the sub-resonance of $H_{I,j}$,
\[
\ln \mu_{b(j)} \leq \sum_k i_k \ln |\lambda_k| \leq \sum_k j_k \ln |\lambda_k|,
\]
where the second inequality holds because $j_k\le i_k$ and $\ln|\lambda_k|<0$ for every $k$.
\end{proof}
We can now explicitly describe the group law of $G$. Let $(\tau_1, h_1)$ and $(\tau_2, h_2)$ be two elements of $G$. We have:
\[
(\tau_1, h_1) \cdot (\tau_2, h_2) = t_{\tau_1} \circ h_1 \circ t_{\tau_2} \circ h_2,
\]
\[
= t_{\tau_1} \circ t_{- \alpha} \circ t_\alpha \circ h_1 \circ t_{\tau_2} \circ h_2,
\]
\[
= t_{\tau_1 - \alpha} \circ h_3,
\]
\[
= (\tau_1 - \alpha, h_3),
\]
where $\alpha = -h_1(\tau_2)$ and $h_3 = t_\alpha \circ h_1 \circ t_{\tau_2} \circ h_2$. By the lemma and the group structure of $SR^*(L)$, $h_3 \in SR^*(L)$.

From the previous calculation, and since $(t_\tau\circ h)^{-1}=h^{-1}\circ t_{-\tau}$, we obtain:
\[
(\tau, h)^{-1} = \big(h^{-1}(-\tau),\ t_{-h^{-1}(-\tau)} \circ h^{-1} \circ t_{-\tau}\big).
\]

We observe that these relations are algebraic in $\tau_1$, $\tau_2$, $h_1$, and $h_2$.
\end{proof}

\end{document}
